% Standalone report fragment: amsmath, amssymb, booktabs, hyperref required.
\section{An optional symmetric-guard and zero-jump composition}
\label{sec:csz}

The companion PA194CSZ has 194 working states. It composes the supplied
symmetric-guard PA194CS machine with the zero-register jump technique.
The exact new native row is
\[
 (27,1):0L40\quad\longmapsto\quad(27,1):1R21.
\]
The complete source still has 3,646 sites. Its only instruction change,
including all successor labels, is
\[
 26:\operatorname{jump}34
 \quad\longmapsto\quad
 26:\operatorname{dec}r_6\,[34,27].
\]
The first successor is the zero branch; a positive decrement lowers the
register and takes the second successor. The new machine's dense table is
\texttt{machine/pa194CSZ.tm}, with initial state zero, working states
$0,\ldots,193$, and separate halt label 194. The tape is initially all zero,
is two-way infinite, and has alphabet $\{0,1\}$; every row moves L or R.

\subsection{A common finite invariant}
For a concrete register vector $v\in\mathbb N^9$, let
$\pi(v)=\sum_{r<9}2^r[v_r>0]$. Abstract states are pairs $(i,m)$ of a source
site and a nine-bit positivity mask. A jump changes only the site.
An increment of $r$ sets bit $r$. A decrement with bit $r$ clear takes its
zero successor with the same mask. With bit $r$ set, it takes its positive
successor with either the same mask or with bit $r$ cleared. Halt has no
successors. The two possible positive-decrement masks overapproximate the
possibilities $v_r>1$ and $v_r=1$.

Induction on concrete steps proves that every concrete successor projects
to one of these abstract successors. The builder computes the least finite
closure from $(25,1)$, corresponding to the full logical seed
$(2,0,0,0,0,0,0,0,0)$. For both the old and new source this closure is the
same set $R$, with
\[
 |R|=263{,}180.
\]
The supplied text certificate encodes a byte array of length $3646\cdot512$,
indexed by $512i+m$. Line $i$ stores the 512 mask-membership bits as a
hexadecimal integer, and the replay checks every abstract edge. At site 26,
\[
 (26,m)\in R\quad\Longrightarrow\quad m\mathbin{\&}64=0.
\]
Moreover the old and new abstract successor lists agree at every element
of $R$. This is checked directly, not inferred merely from the sizes of
the two closures.

Consequently $r_6=0$ whenever a concrete source run from the stated seed
visits site 26. The old jump and the new decrement then have the identical
successor, with every register unchanged. At every other site the source
instructions coincide. Induction on the source-step count proves equality
of the complete old and new source executions, including halting.

The invariant has been recomputed after the symmetric-guard change.
The older cached-controller zero-jump proof had a different closure and
cannot by itself justify this composition. The finite source invariant,
its soundness induction, and the resulting trace transformation are all
formalizable in PA.

\subsection{Binding the new literal table}
The construction checker verifies 4,746 arbitrary-register source
outcomes. These include both branches of the new decrement, although the
positive branch is excluded from the actual logical run by the invariant.
It uses 1,513,882 symbolic steps, with maximum 508 per outcome.
An independent literal-tape checker, using no arithmetic-box oracle,
rederives four generic list traversals from the exact table and verifies
all 4,746 outcomes in 1,100,566 symbolic steps, with maximum 403. The
smallest positive-duration constant is three. These symbolic counts
measure finite proof-case work, not bounds on unbounded numerical runs.
The quantified frames contain arbitrary natural register values, an
affine gap, and an arbitrary finite nonempty spare list.

Separate literal simulations verify that the new machine reaches the
entire clean source frame at site 25 after exactly 58,014 blank-tape
steps, including all markers, register blocks, spare blocks, head
position and zero exterior. This differs from PA194CS's 57,768-step
startup; no runtime improvement is claimed. The dense and native tables
are checked by the explicit 195-element control bijection, including halt.

The source equivalence above, the exact physical simulations and startup,
and the explicitly inherited PA194CS ordinary consistency calibration
therefore transfer its PA-provable nonhalting equivalence to PA194CSZ.
This is an ordinary mathematical proof with executable certificates;
it is not a new end-to-end Lean or Rocq theorem. The contribution is the
verified composition and the newly computed source invariant, not a
state-count improvement.
